\runninghead{Each Element in Its Setting}
\section{Each Element in Its Setting}
\label{sec:each-element}

Each appointed text is read here first in its own place: the chapter or psalm
it comes from, the words the Missal keeps and leaves, and what its principal
witnesses say of it. The three readings that follow take these findings up
across the whole Mass.

\subsection{Introit: the Lamb before the throne}
\label{sec:element-introit}

The antiphon comes from the throne vision of Apoc 4--5. In the right hand of
him who sits on the throne John sees a book sealed with seven seals, and no one
in heaven or on earth is found worthy to open it. One of the ancients tells him
not to weep, for ``the lion of the tribe of Juda'' has prevailed; John looks
and sees ``a Lamb standing as it were slain'' (5:5--6). The four living
creatures and the ancients then sing a new canticle to the Lamb who was slain
and has redeemed men to God ``out of every tribe and tongue and people and
nation'' and ``made us to our God a kingdom and priests, and we shall reign on
the earth'' (5:9--10). Only after that do the angels, thousands of thousands,
take up the sevenfold acclamation of verse~12 that the Introit sings. The
Introit's second sentence comes from John's greeting at the head of the book,
from Jesus Christ ``the faithful witness, the first begotten of the dead and
the prince of the kings of the earth, who hath loved us and washed us from our
sins in his own blood'' (1:5).

Victorinus of Pettau reads the opening of the book as the overcoming of death:
none was found worthy save the Lamb, and because he overcame death as
testator it was just that he be made heir and possess the human members (on
5:2--5). St Bede reads the Lion and the Lamb as one Lord, a Lamb in dying
innocently and a Lion in conquering death (\work{Expl.\ Apoc.}\ I, on 5:5--6).
Blessed Ildefonso Schuster, commenting on this Mass's Introit in 1930, notes
that John recalls the Passion, the Lamb ``immolated and slain,'' before he
shows the Lamb's glory, because ``the exaltation of Jesus Christ is, according
to St Paul, in proportion to the humiliation he endured'' (\work{The
Sacramentary} V, p.~201). The verb that the antiphon sets at its centre,
\latin{accipere}, and its strange object, \latin{divinitatem}, receive their
fullest treatment in the first reading below.

The verse is the opening prayer of Psalm~71, \latin{Deus, iudicium tuum Regi
da: et iustitiam tuam Filio Regis}. The psalm is titled \latin{in Salomonem}
and asks for the king a reign of justice to the poor, a peace that lasts as
long as the moon, dominion from sea to sea, the homage of the kings of
Tharsis, Arabia and Saba, and the blessing of all nations in him; it closes
\latin{Defecerunt laudes David, filii Jesse} (71:20). Its witnesses differ on
its literal sense. St Augustine holds that the title is figurative, because the
psalm says what could not be true of Solomon according to the flesh: Solomon
means peacemaker, the true Peacemaker is the Mediator, and the King who is
also the King's Son is Christ, to whom the Father has given all judgment (Jn
5:22; \work{En.\ in Ps.}\ 71.1). Cassiodorus and Theodoret of Cyrus read it the
same way, Theodoret adding that Solomon neither ruled the ends of the earth nor
ended his life in glory. St Robert Bellarmine gives the literal sense as
David's prayer for his son Solomon, ``just anointed king,'' who asks not riches
but the grace of judging rightly, ``for kings were created for the people, not
the people for kings''; he allows the reading of Christ beside it, and says
the psalm ``now begins to pass from Solomon to Christ'' only at the verse ``he
shall continue with the sun'' (\work{Explanatio}, Ps~71). Schuster calls the
verse ``distinctly Messianic,'' the Messianic office conferred upon Jesus
(pp.~201--202).

\subsection{Collect: all things restored in the Son}
\label{sec:element-collect}

The Collect was composed for the feast and first printed in 1925. Its
principal clause, \latin{qui in dilecto Filio tuo \dots\ omnia instaurare
voluisti}, takes its verb from the Vulgate of Eph 1:10, \latin{instaurare omnia
in Christo}, the ``re-establishing'' or summing up of all things in Christ, and
its \latin{dilecto Filio} from the ``beloved Son'' of Eph 1:6 and the ``Son of
his love'' of Col 1:13. Two witnesses tie that verse of Ephesians to the
Epistle. St John Chrysostom, on Col 1:20, says the things in heaven are
reconciled because Christ brought man up to the angels, and joins this to
``to sum up all things in Christ'' (\work{Hom.\ in Col.}\ 3). Theodoret reads
the \latin{in ipsum} of Col 1:16 of the work done through the Incarnation and
cites \latin{Instaurare omnia in Christo} for it (\work{In Col.}\ 1). The
title \latin{universorum Rege} names what Col 1:16--17 states of the Son, in
whom all things were created and hold together, and what the Preface will
call him.

The petition gives the condition of the nations and the remedy. They are
\latin{familiae gentium}, the families of the peoples, and they are
\latin{peccati vulnere disgregatae}: their scattering is a wound, and its cause
is sin. Schuster reads the prayer as inspired by St Paul: sin is ``a
disruptive force'' that ruined the unity of the human race, which Christ the
new Adam restores (p.~202). What the Collect asks for them is subjection, but
to a rule whose quality it names first, \latin{suavissimo imperio}, ``most
gentle'' or, in the Lasance, ``most loving.'' How a kingdom not of this world
can ask the subjection of nations is the question of the second reading.

\subsection{Epistle: translated into the kingdom of the Son}
\label{sec:element-epistle}

Paul writes to the Colossians in bonds (Col 4:3, 18), and the passage stands
inside his prayer for them (1:9--11) and before its application to them as
once enemies, now reconciled (1:21--23). Its first half is a thanksgiving to
the Father, who has made us worthy of ``the lot of the saints in light,''
delivered us ``from the power of darkness,'' and ``translated us into the
kingdom of the Son of his love'' (1:12--13). The 1962 page sets the words of
verse~13 as the caption of the engraved headpiece above the formulary. The
second half is the hymn to the Son as image of the invisible God and firstborn
of every creature, in whom, through whom and unto whom all things were made,
visible and invisible, thrones and dominations and principalities and powers;
as head of the body, the Church, and firstborn from the dead; as the one in
whom the fulness dwells, through whom the Father reconciles all things,
``making peace through the blood of his cross'' (1:15--20).

Chrysostom stresses the verb of verse~13: God did not only deliver us from
darkness but translated us into a kingdom, and called it the kingdom of the
Son of his love, giving us the same honour as the Son (\work{Hom.\ in
Col.}\ 2). In the next homily he insists, on verse~15, that ``firstborn'' does not mean ``first
created,'' and on verse~18 that ``firstborn from the dead'' means the
firstfruits of the resurrection. St Thomas Aquinas glosses \latin{transtulit}
as \latin{ut essemus regnum Dei}, ``that we might be the kingdom of God,'' and
joins to it the Gospel's \latin{regnum meum non est de hoc mundo} and the
Apocalypse's ``thou hast made us to our God a kingdom'' (5:10); he assigns the
redemption to Christ \latin{inquantum homo, factus \dots\ pro nobis
sacrificium}, and the remission of sins to him as God (\work{Super Col.}\ 1
lect.~3). On verse~18 he makes Christ head of the Church \latin{secundum quod
huiusmodi}, as man, first in grace and first in glory, and calls the
resurrection a second generation (lect.~5). Theodoret marks the same turn at
verse~18, where Paul passes, he says, from theology to the economy, and Christ
is our head according to the humanity; at verse~20 he says Christ made our
reconciliation by suffering and offering the sacrifice for us, and joined the
heavenly to the earthly. Schuster sees in the Epistle Christ's primacy over all
creation, the angels included, founded on the hypostatic union, and his
headship of the Church (p.~202).

The Missal's text departs from the Vulgate only at its frame: it opens
\latin{Fratres: Gratias agimus} and closes \latin{in Christo Iesu Domino
nostro}.

\subsection{Gradual: from sea to sea}
\label{sec:element-gradual}

The Gradual takes Ps 71:8 and 71:11 and leaves out the verses between them,
in which the Ethiopians fall down before the king, his enemies lick the
ground, and the kings of Tharsis, the islands, Arabia and Saba bring gifts.
Its respond is the extent of the dominion, its verse the homage of all kings
and the service of all nations.

The Fathers read the geography of verse~8 as the whole earth. For St Augustine
``from sea to sea'' is from any end of the earth to any other, the Church
spread through the whole world, and ``from the river'' is the Jordan, where
the Father's voice was heard at the baptism and from which the Gospel of the
kingdom went out to all nations (\work{En.\ in Ps.}\ 71.11). Cassiodorus adds
that the verse cannot belong to Solomon, who reigned in Judaea alone, and that
from the Jordan \latin{per totum orbem terrarum saluberrima regula sacri
baptismatis emanavit}, the saving rule of holy baptism flowed out through the
whole world. Theodoret takes the seas as the outer seas that bound the
inhabited world, and the river as the Jordan, from which the preaching ``Do
penance, for the kingdom of heaven is at hand'' began. Bellarmine, following
Augustine, prefers the Jordan, where the kingdom had its rise.

Schuster reads the Gradual's prophecy as being fulfilled while the Church
extends her peaceful conquests, and entire only at the judgment (p.~202; his
words are given in the second reading). He also treats the Gradual and the
Alleluia together as ``the double chant after the Epistle,'' which sets out a
threefold
catholicity: of place, \latin{a mari usque ad mare}; of subjects, \latin{omnes
reges \dots\ omnes gentes}, kings and ``wretched pariahs'' alike; and of ages,
\latin{potestas aeterna \dots\ non auferetur} (p.~203). That is the one join
between two elements of the Mass that he makes himself. What the kings' homage
means, and when it is given, is developed in the second reading.

\subsection{Alleluia: an everlasting power}
\label{sec:element-alleluia}

Daniel dates his vision to the first year of Baltasar, king of Babylon (Dan
7:1). Four beasts rise from the sea; thrones are set, the Ancient of days sits,
the books are opened and the beast is slain. Then ``one like the Son of man
came with the clouds of heaven, and he came even to the ancient of days: and
they presented him before him'' (7:13), and he is given power and glory and a
kingdom that all peoples serve, an everlasting power and an indestructible
kingdom (7:14). In the interpretation, ``the saints of the most high God''
receive the kingdom and possess it for ever (7:18, 27). The Alleluia sings only
the second half of verse~14: the power as possessed and lasting, not the
giving of it.

St Jerome identifies the son of man coming with the clouds with the stone of
Nabuchodonosor's dream, and reads both of the Son of God's
assumption of human flesh. His reading of the giving of power by the
Apostle's \latin{Qui cum in forma Dei esset} (Phil 2:6--8) belongs to the first
reading below, and his argument against Porphyry, that no man can hold a power
which no end encloses, to the third (\work{In Dan.}\ II, on 7:13--14). On
7:17--18 he holds that the saints receive a heavenly, not an earthly, kingdom:
\latin{cesset ergo mille annorum fabula}, let the tale of the thousand years
cease. The scholia on Daniel transmitted under the name of Hippolytus say that
the Father put all things in subjection to the Son and showed him
first-begotten of God, of a virgin and from the dead: King of all in heaven
because the Word born before all, of all on earth because made man, of all
under the earth because numbered among the dead. Aquinas cites this verse in
his exposition of the Gospel, for the origin and the permanence of Christ's
royal power (\work{Super Io.}\ 18 lect.~6). Schuster sees its partial
fulfilment in the Church's two thousand years against the empires that have
passed away (pp.~202--203).

\subsection{Gospel: the king before the governor}
\label{sec:element-gospel}

St John places the scene in the praetorium, the governor's hall, early on the
morning before the Pasch; the accusers will not enter it, so that Pilate goes
in and out between them and the prisoner (Jn 18:28--29). The Gospel is his
first questioning of Jesus inside. What follows it in John's account is
Pilate's ``What is truth?'' (18:38), the choice of Barabbas, the scourging, the
crown of thorns and the purple, \latin{Ecce homo}, the chief priests' ``We have
no king but Caesar,'' and the handing over to be crucified (18:39--19:16).

Each phrase of the dialogue has its own difficulty. Jesus answers Pilate's
question with a question; St Augustine reads it as exposing the vain notion of
his kingdom that Gentiles and Jews alike had formed (\work{Tract.\ in Io.}\
115.1). Verse~36 denies that the kingdom belongs to this world, and the Latin
words the denial three times: the kingdom is not \latin{de hoc mundo}; were it
\latin{ex hoc mundo}, his servants would strive for him; and now it is not
\latin{hinc}. The Douay renders the first two alike, ``of this world,'' and the
third ``from hence.'' Augustine draws from the words the distinction of
``of'' this world from ``in'' it, and of ``from hence'' from ``here'' (115.2),
which the second reading takes up. ``Thou sayest that I am a king'' he reads as
neither a denial nor a confession of the king Pilate imagines (115.3). ``For
this was I born'' divides the witnesses. Augustine reads it of the temporal
nativity, ``by being born of the Virgin'' (115.4); Chrysostom of the eternal
generation, for he who was born a king has all his attributes by generation (\work{Hom.\ in Io.}\ 83.4); Aquinas sets the
two readings side by side, the second \latin{sicut Deus ex Deo, ita rex ex
rege}, and adds that ``for this came I into the world'' is the temporal coming
even on Chrysostom's reading (\work{Super Io.}\ 18 lect.~6). St Cyril of
Alexandria reads ``came into the world'' as ``when I became Man'' (\work{In
Io.}\ XII). The whole dialogue, and the denial at its centre, is the text of the
second reading.

\subsection{Offertory: ask of me}
\label{sec:element-offertory}

The antiphon is one verse of the second psalm, which Acts 4:25 ascribes to
David. The psalm opens with the nations raging and the kings of the earth
standing up ``against the Lord and against his Christ'' (2:1--2); he who dwells
in heaven laughs at them; the king speaks: ``I am appointed king by him over
Sion, his holy mountain,'' and recites the Lord's decree, ``Thou art my son,
this day have I begotten thee'' (2:6--7). The Offertory's verse follows at once.
After it the psalm gives the king a rod of iron (2:9) and turns to the rulers:
``And now, O ye kings, understand: receive instruction, you that judge the
earth. Serve ye the Lord with fear: and rejoice unto him with trembling''
(2:10--11).

The Missal sings only verse~8, but its ``Ask of me'' is spoken to the Son of
verse~7, and the Fathers read the two verses together. Their reading of that
pair, of what it means for the eternal Son to ask, is the centre of the first
reading; their reading of the kings of verses 10--11 belongs to the second.
Schuster, on this antiphon, says that to ``the nations and the kings who rise
up against'' God's design Christ shows Ps 2:8 as ``the solemn deed of
investiture which he has received from his Father'' (p.~203).

\subsection{Secret: the victim of reconciliation}
\label{sec:element-secret}

The Secret, composed for the feast, names the sacrifice by its effect:
\latin{hostiam \dots\ humanae reconciliationis}, the victim by which mankind
is reconciled, in the vocabulary of Col 1:20 and of 2 Cor 5:18--19, ``God
\dots\ hath reconciled us to himself by Christ.'' Its relative clause
identifies the victim: \latin{quem sacrificiis praesentibus immolamus}, the
Son himself, immolated in the sacrifice now offered. And it asks of him the
gifts that are the opposite of the Collect's wound: for all nations,
\latin{unitatis et pacis dona}, the gifts of unity and peace. Schuster gives
the Secret's text without comment (p.~204). After it the Missal directs the
proper Preface.

\subsection{Preface: priest and king}
\label{sec:element-preface}

The Preface, too, was composed for the feast and first printed with the Mass
in 1925, in place between the Secret and the Communion; the 1962 Missal prints
it among the Prefaces and appoints it for the feast and for the votive Mass of
Christ the King (RGMR 492). Its argument runs in one sentence. The Father
anointed his only-begotten Son \latin{oleo exsultationis}, the ``oil of
gladness'' of Ps 44:8, which Heb 1:9 applies to the Son, and anointed him
both \latin{Sacerdotem aeternum} and \latin{universorum Regem}. The purpose of
the anointing is double and ordered: first, that offering himself on the altar
of the cross as a spotless and peace-making victim he might accomplish
\latin{redemptionis humanae sacramenta}; then, all creatures being subjected
to his rule, that he might hand over to the Father's majesty an eternal and
universal kingdom. The verb \latin{traderet} is the ``delivered up'' of 1 Cor
15:24, ``when he shall have delivered up the kingdom to God and the Father.''
The kingdom is then named by its goods, in three groups: \latin{veritatis et
vitae}; \latin{sanctitatis et gratiae}; \latin{iustitiae, amoris et pacis}.

Schuster's commentary treats the Preface as a text. Its conception of Christ as Pontiff and King ``who reconquers and restores to his
Heavenly Father the realm of creation \dots\ that God might eternally be all in
all'' is, he writes, ``borrowed from St Paul,'' and its compiler gave it the
lyrical character of the early eucharistic \latin{Hymnus} (p.~204). Bede, on
the clause of Apoc 1:6 that the Introit leaves out, gives the same joining of
king and priest by self-offering: \latin{Quia Rex regum et Sacerdos coelestis
se offerendo pro nobis nos suo corpori adunavit, nemo sanctorum est qui
spiritualiter sacerdotii officio careat}, because the King of kings and
heavenly Priest, by offering himself for us, has united us to his body, there
is no saint who lacks the office of the priesthood in the spirit (\work{Expl.\
Apoc.}\ I, on 1:6). \work{Quas primas} joins the kingly office to the offices
of redeemer and priest (n.~16). The handing over of the kingdom, and its
relation to a reign that never ends, is the text of the third reading.

\subsection{Communion: the Lord shall sit King for ever}
\label{sec:element-communion}

Psalm~28 is titled, in the Douay, ``A psalm for David, at the finishing of the
tabernacle.'' It summons the sons of God to give glory to the Lord, and then
hears his voice seven times: upon the waters, in power, in magnificence,
breaking the cedars of Libanus, dividing the flame of fire, shaking the
desert, preparing the stags and discovering the thick woods, while in his
temple all speak his glory. It ends with the two verses from which the
Communion takes its halves: ``The Lord maketh the flood to dwell: and the Lord
shall sit king for ever. The Lord will give strength to his people: the Lord
will bless his people with peace.'' The antiphon keeps the session and the
blessing and leaves out the flood and the strength.

The witnesses do not agree on what the flood is, and so they place the
session differently. St Thomas Aquinas gives first the plain sense of the
Hebrew, \latin{Dominus diluvio sedit}: his judgments were manifest from the
beginning, as when he brought the deluge, and he will sit King for ever,
judging the peoples in equity. Then he gives three mystical readings: the saints dwelling secure in the
Church as in Noe's ark amid the flood of tribulations; the world and its carnal
men, the \latin{diluvium}, brought to dwell in his temple when they are
converted, \latin{et sedebit rex in aeternum}; and the baptismal waters, where
he dwells by the effect of grace (\work{Postilla super Psalmos} 28). St
Augustine reads the flood as this world, in which the Lord dwells first in his
saints, kept in the Church as in the ark (\work{En.\ in Ps.}\ 28.10). St Basil
the Great calls the grace of baptism a deluge, and Cassiodorus sees the Holy
Trinity dwelling in the waters of baptism; Theodoret, the world drowned by the
torrent of lawlessness and made a new creation. Bellarmine reports both the
deluge and baptism and will not contradict those who hold them, ``being of
great weight and high position,'' while his own second reading places the
flood at the day of judgment (\work{Explanatio}, Ps~28). Schuster reads the
antiphon of the communicant, in whose heart the King is enthroned (p.~204).
The third reading follows these witnesses into the session for ever.

\subsection{Postcommunion: under the standards of the King}
\label{sec:element-postcommunion}

The Postcommunion begins from what has just been received, \latin{immortalitatis
alimoniam}, the food of immortality, and asks a share in the reign: that those
who glory to fight under the standards of Christ the King, \latin{sub Christi
Regis vexillis militare}, may reign with him for ever \latin{in caelesti sede}.
The last clause echoes 2 Tim 2:12, ``If we suffer, we shall also reign with
him,'' which Theodoret cites on Col 1:13, and the Apocalypse's ``we shall
reign on the earth'' (5:10). Its figure makes the faithful soldiers, as no
other text of the Mass does, and it stands a few minutes after a Gospel in
which the King says that his servants do not fight. Schuster's note on the prayer says what the
standard is: ``The Royal Standard of Jesus Christ is the Cross: Dicite in
nationibus: regnavit a ligno Deus; for the kingdom of Christ is that of
obedience, humility and sacrifice'' (p.~204). The second reading takes up the
tension between the standard and the Gospel.
